\chapter{SYSTEM IMPLEMENTATION AND TESTING}
\label{chap:implementation}

\section{System Implementation}

The Implementation phase for the Store Management System of BUBT involves the transformation of system design into a functional, secure, and institutional-grade platform. This phase ensures that the features and automation logic are correctly implemented and validated according to university requirements.

The implementation follows a modular approach, with each system component developed, tested, and integrated independently before final system testing. This approach minimizes integration issues and allows for parallel development of different modules.

\subsection{Technology Stack}

The system is built using the following technologies:

\begin{itemize}
    \item \textbf{Backend Framework}: Laravel 12.0 - A PHP framework for web application development
    \item \textbf{Programming Language}: PHP 8.2 - Server-side programming language
    \item \textbf{Database}: MySQL - Relational database management system
    \item \textbf{Frontend Framework}: Tailwind CSS - Utility-first CSS framework
    \item \textbf{JavaScript Library}: Alpine.js - Lightweight JavaScript framework for interactivity
    \item \textbf{Build Tool}: Vite - Next-generation frontend build tool
    \item \textbf{Authentication}: Laravel Breeze - Minimal authentication implementation
    \item \textbf{Authorization}: Spatie Permission - Role and permission management
    \item \textbf{PDF Generation}: Barryvdh Laravel DomPDF - PDF document generation
    \item \textbf{Image Processing}: Intervention Image - Image handling library
    \item \textbf{Real-time Messaging}: Chatify with Pusher - Real-time chat functionality
\end{itemize}

This technology stack provides a secure, scalable, and maintainable foundation for the inventory management system. All components are open-source and widely supported by active communities.

\subsection{Frontend Implementation}

The frontend implementation focuses on creating a responsive, user-friendly interface that supports all system functionalities. Key aspects include:

\begin{enumerate}
    \item \textbf{Authentication Interface}
    \begin{itemize}
        \item Login page with email and password authentication
        \item Registration form with role selection
        \item Password recovery with email verification
    \end{itemize}

    \item \textbf{Dashboard}
    \begin{itemize}
        \item Summary cards showing key statistics
        \item Navigation sidebar with module access
        \item Quick action buttons
        \item User profile and notification area
    \end{itemize}

    \item \textbf{Data Tables}
    \begin{itemize}
        \item Sortable and searchable tables
        \item Pagination controls
        \item Action buttons for CRUD operations
    \end{itemize}

    \item \textbf{Forms}
    \begin{itemize}
        \item Input validation with error messages
        \item Date pickers for date selection
        \item File upload for images and documents
        \item Dynamic form fields for variable data
    \end{itemize}
\end{enumerate}

The user interface follows consistent design patterns throughout the application, reducing the learning curve for users. Tailwind CSS ensures responsive design that adapts to different screen sizes.

\begin{figure}[H]
    \centering
    \includegraphics[width=0.9\textwidth]{figure/dashboard.png}
    \caption{System Dashboard Interface}
    \label{fig:dashboard}
\end{figure}

\subsection{Backend Implementation}

The backend implementation handles all business logic, database operations, and system security. Key components include:

\begin{enumerate}
    \item \textbf{Controllers}
    \begin{itemize}
        \item Authentication controllers for login, registration, and password management
        \item Product, Category, Brand, and Supplier controllers for entity management
        \item Purchase, Requisition, and Issue controllers for workflow management
        \item Report controller for generating various reports
        \item Role and Permission controller for access control
    \end{itemize}

    \item \textbf{Models}
    \begin{itemize}
        \item Eloquent models for all database entities
        \item Relationships defined for data navigation
        \item Accessors and mutators for data formatting
    \end{itemize}

    \item \textbf{Middleware}
    \begin{itemize}
        \item Authentication middleware for protected routes
        \item Role-based permission checking middleware
        \item Custom middleware for specific requirements
    \end{itemize}

    \item \textbf{Services}
    \begin{itemize}
        \item Business logic extraction into service classes
        \item Report generation services
        \item PDF generation services
    \end{itemize}
\end{enumerate}

The backend follows Laravel's MVC architecture, ensuring separation of concerns and maintainable code. Routes are organized by module, making it easy to locate and modify endpoint definitions.

\section{System Modules}

The system consists of several integrated modules that work together to provide comprehensive inventory management capabilities.

\subsection{Authentication and Authorization Module}

This module handles user identity verification and access control:

\begin{itemize}
    \item User registration with email verification
    \item Secure login with session management
    \item Password reset functionality
    \item Role assignment and permission management
    \item Department and semester management
\end{itemize}

The module uses Laravel Breeze for authentication scaffolding and Spatie Permission for role-based access control. Users are assigned roles that determine their access to different system functions.

\begin{figure}[H]
    \centering
    \includegraphics[width=0.9\textwidth]{figure/login_page.png}
    \caption{Login Page Interface}
    \label{fig:login}
\end{figure}

\subsection{Product Management Module}

This module provides comprehensive product management:

\begin{itemize}
    \item Product creation with image upload
    \item Category and subcategory assignment
    \item Brand association
    \item SKU generation
    \item Stock quantity tracking
    \item Role-based product visibility
    \item Product search and filtering
\end{itemize}

The module integrates with Intervention Image for handling product images, including automatic resizing and optimization.

\subsection{Supplier Management Module}

This module manages vendor information:

\begin{itemize}
    \item Supplier profile creation
    \item Contact information management
    \item Purchase history tracking
    \item Supplier status management
\end{itemize}

The supplier database enables efficient vendor management and supports purchase order creation.

\subsection{Purchase Management Module}

This module handles the complete purchase workflow:

\begin{itemize}
    \item Purchase order creation with multiple items
    \item Supplier selection
    \item Semester and department association
    \item Tracking number and reference number
    \item File attachment for supporting documents
    \item Automatic stock update on receiving
    \item Purchase invoice generation
\end{itemize}

The module maintains complete audit trail of all purchase transactions, supporting accountability and reporting requirements.

\begin{figure}[H]
    \centering
    \includegraphics[width=0.9\textwidth]{figure/purchase_management.png}
    \caption{Purchase Management Interface}
    \label{fig:purchase}
\end{figure}

\subsection{Requisition Management Module}

This module handles product requests from departments:

\begin{itemize}
    \item Requisition creation by department users
    \item Item selection with quantity specification
    \item Status tracking (Pending, Approved, Rejected)
    \item Approval workflow for administrators
    \item Requisition history viewing
\end{itemize}

The workflow ensures proper authorization of product requests while maintaining efficiency in processing.

\subsection{Issue Management Module}

This module manages product issuance:

\begin{itemize}
    \item Issue creation from approved requisitions
    \item Automatic stock deduction
    \item Issue date and reference tracking
    \item Issue invoice generation
    \item Issue history reporting
\end{itemize}

The module integrates with requisition management to ensure that only approved requests are fulfilled.

\subsection{Damage Product Management Module}

This module tracks damaged items:

\begin{itemize}
    \item Damage product reporting
    \item Item and quantity documentation
    \item Loss amount calculation
    \item Reason tracking
    \item Damage report generation
\end{itemize}

Proper damage tracking enables accurate inventory records and supports insurance and audit requirements.

\subsection{Return Purchase Management Module}

This module handles product returns:

\begin{itemize}
    \item Return purchase order creation
    \item Item and quantity specification
    \item Return reason tracking
    \item Automatic stock adjustment
    \item Return invoice generation
\end{itemize}

The module maintains accurate records of returned items for accounting and inventory purposes.

\subsection{Issue Return Management Module}

This module handles returned issued items:

\begin{itemize}
    \item Issue return creation
    \item Return reason documentation
    \item Stock restoration on return
    \item Return tracking and reporting
\end{itemize}

This feature ensures flexibility in handling situations where issued products need to be returned to inventory.

\subsection{Quotation Management Module}

This module supports procurement planning:

\begin{itemize}
    \item Quotation creation from suppliers
    \item Item and price specification
    \item Validity period tracking
    \item Quotation comparison
    \item Quotation approval workflow
\end{itemize}

The quotation module facilitates competitive procurement by enabling comparison of supplier quotes.

\subsection{Report Generation Module}

This module provides comprehensive reporting:

\begin{itemize}
    \item Stock status reports
    \item Purchase reports (with date range filtering)
    \item Purchase return reports
    \item Sale reports
    \item Sale return reports
    \item Damage product reports
    \item Product lifetime reports
    \item PDF export for all reports
\end{itemize}

Reports are generated using Barryvdh Laravel DomPDF, providing professional-looking documents that can be easily shared and archived.

\begin{figure}[H]
    \centering
    \includegraphics[width=0.9\textwidth]{figure/report_generation.png}
    \caption{Report Generation Interface}
    \label{fig:report}
\end{figure}

\subsection{Chat/Messaging Module}

This module provides real-time communication:

\begin{itemize}
    \item User-to-user messaging
    \item Message history
    \item Online status indication
    \item File sharing capability
\end{itemize}

The Chatify integration with Pusher provides real-time messaging functionality for improved coordination between administrators and department users.

\section{System Setup}

The system setup involves configuring the software environment, database, and security infrastructure.

\subsection{Web Server Configuration}

The application is deployed on a web server supporting PHP and MySQL:

\begin{itemize}
    \item Apache or Nginx web server configuration
    \item PHP 8.2 runtime configuration
    \item Database connection setup in .env file
    \item Storage and cache directory permissions
\end{itemize}

The configuration ensures proper handling of HTTP requests and background processing.

\subsection{Database Setup}

The database setup includes:

\begin{itemize}
    \item MySQL database creation
    \item Migration execution for schema creation
    \item Seeding for initial data population
    \item Connection configuration
\end{itemize}

Database migrations ensure consistent schema across development and production environments.

\section{Testing Strategies}

Testing is conducted to ensure that the system remains efficient, automated, and secure from unauthorized access.

\subsection{Unit Testing}

Unit testing validates individual components in isolation:

\begin{itemize}
    \item Model functionality testing
    \item Controller method validation
    \item Service class verification
    \item Helper function testing
\end{itemize}

Unit tests identify bugs early in the development process and ensure that individual components function correctly.

\subsection{Integration Testing}

Integration testing verifies that components work together:

\begin{itemize}
    \item Database query validation
    \item Controller-model interaction
    \item Middleware behavior
    \item Route handling
\end{itemize}

Integration tests ensure that the system functions correctly as a whole.

\subsection{Functional Testing}

Functional testing validates complete user workflows:

\begin{itemize}
    \item Authentication flow testing
    \item CRUD operation validation
    \item Report generation verification
    \item Workflow completion testing
\end{itemize}

Functional tests ensure that the system meets user requirements and functions as specified.

\subsection{Security Testing}

Security testing ensures protection against vulnerabilities:

\begin{itemize}
    \item Authentication bypass attempts
    \item SQL injection prevention
    \item Cross-site scripting (XSS) protection
    \item CSRF token validation
    \item Session hijacking prevention
\end{itemize}

Security testing is critical for a system that handles sensitive institutional data.

\subsection{User Acceptance Testing}

User acceptance testing involves end users evaluating the system:

\begin{itemize}
    \item Usability assessment
    \item Functionality verification
    \item Workflow validation
    \item User feedback collection
\end{itemize}

UAT ensures that the system meets user expectations and is ready for production deployment.

\section{Testing and Quality Assurance}

The Quality Assurance (QA) process focuses on validating system features while ensuring the system minimizes manual errors in university resource management.

\subsection{Test Objectives}

The primary testing objectives include:

\begin{enumerate}
    \item \textbf{Security Validation}: Ensuring that the login page prevents unauthorized access to sensitive university resources.

    \item \textbf{Reliability}: Verifying that the system provides fast and reliable data storage and retrieval while maintaining institutional data integrity.

    \item \textbf{Operational Continuity}: Confirming that self-service features support uninterrupted store management operations.

    \item \textbf{Data Accuracy}: Ensuring that all calculations and stock updates are accurate and consistent.
\end{enumerate}

\subsection{Quality Assurance Standards}

The QA process ensures that the final product meets academic, technical, and professional standards:

\begin{itemize}
    \item Code quality review and adherence to coding standards
    \item Security vulnerability assessment and remediation
    \item Performance optimization and load testing
    \item Cross-browser and cross-device compatibility verification
    \item Documentation completeness and accuracy
\end{itemize}

Through automation, security validation, and error reduction mechanisms, the system is verified as a reliable solution for long-term university resource management.

\section{Summary}

The Implementation and Testing phase successfully transformed the system design into a functional and secure platform. The modular implementation approach ensured that all components work together effectively while maintaining code quality and security standards.

The comprehensive testing strategy validated that the system meets all functional requirements while providing the security and reliability expected of an institutional application. User acceptance testing confirmed that the system is ready for production deployment.